% GENERATED FILE. Edit canonical lessons, then run npm run generate:books.
% canonical-source-hash: fnv1a64:c35b81b41c864720
% canonical-lessons: TE-C518-gayapadu, TE-C518-boppi, TE-C518-tulu, TE-C518-kuntu, TE-C518-raktasravam

\chapter{Slips and Falls}
\label{ch:te-slips-and-falls}
\hlchaptermodality{518}

\begin{chapteropening}
\textbf{By the end of this chapter:} \emph{I can say that someone got hurt, got a bump, staggered, limped or bled.}

Complete the last lesson of chapter 518: I can say that someone got hurt, got a bump, staggered, limped or bled.
\end{chapteropening}

\section[\texorpdfstring{\te{గాయపడు} (gāyapaḍu)}{గాయపడు (gāyapaḍu)}]{\te{గాయపడు} (gāyapaḍu) — to get hurt; to be injured}

\label{lesson:TE-C518-gayapadu}



\begin{quote}
Before the new one: say the Telugu for to throw up, then the Telugu for joint pains.
\end{quote}

\subsection*{You'll want to know: \te{గాయపడు}}
\textbf{\te{గాయపడు}} — \emph{gāyapaḍu} — \textquotedblleft{}to get hurt; to be injured\textquotedblright{}.

\begin{quote}
\te{ఆ} \te{ప్రమాదంలో} \te{ఇద్దరు} \te{గాయపడ్డారు}.
\end{quote}

\begin{quote}
\emph{ā pramādaṁlō iddaru gāyapaḍḍāru.}
\end{quote}

\begin{quote}
Two people were injured in that accident.
\end{quote}

Past (I …): \te{గాయపడ్డాను}, \emph{gāyapaḍḍānu}.

\te{గాయం}, a wound, + \te{పడు}. Formal; the news says it.

\begin{cousinweb}
\textbf{\te{గాయపడు}} (\emph{gāyapaḍu}) joins \te{గాయం}, \textquotedblleft{}a wound\textquotedblright{}, usually traced to Sanskrit \emph{ghāta}, \textquotedblleft{}killing; a blow, a bruise\textquotedblright{} (CDIAL 4460), to the native verb \emph{paḍu}, \textquotedblleft{}to feel, to suffer\textquotedblright{} (DEDR 3853).
\end{cousinweb}

\begin{grammarlens}[title={What you've built}]
One more word for describing what you lived and felt.
\end{grammarlens}

\subsection*{Your turn}
\begin{itemize}
\raggedright
  \item \emph{Say it:} \emph{gāyapaḍu}
  \item \emph{Say it:} \emph{gāyapaḍu}, once more
  \item \emph{Say it:} \emph{ā pramādaṁlō iddaru gāyapaḍḍāru.} — Two people were injured in that accident.
  \item \emph{Recall:} say the Telugu for to throw up, then the Telugu for joint pains, then say \emph{gāyapaḍu} again
\end{itemize}

\begin{tcolorbox}[breakable,colback=teal!4,colframe=teal!35!black,title={Before you move on}]
What does \te{గాయపడు} mean? (\textquotedblleft{}To get hurt; to be injured\textquotedblright{}.) Say it once more.
\end{tcolorbox}
\section[\texorpdfstring{\te{బొప్పి} (boppi)}{బొప్పి (boppi)}]{\te{బొప్పి} (boppi) — a bump; a lump on the head}

\label{lesson:TE-C518-boppi}



\begin{quote}
Before the new one: say the Telugu for joint pains, then the Telugu for to get hurt.
\end{quote}

\subsection*{You'll want to know: \te{బొప్పి}}
\textbf{\te{బొప్పి}} — \emph{boppi} — \textquotedblleft{}a bump; a lump on the head\textquotedblright{}.

\begin{quote}
\te{తలుపుకు} \te{తగిలి} \te{తమ్ముడి} \te{నుదుటి} \te{మీద} \te{బొప్పి} \te{కట్టింది}.
\end{quote}

\begin{quote}
\emph{talupuku tagili tammuḍi nuduṭi mīda boppi kaṭṭindi.}
\end{quote}

\begin{quote}
My little brother hit the door and a bump came up on his forehead.
\end{quote}

\te{బొప్పి} \te{కట్టు}, a bump comes up. It belongs with \te{పొంగు}, to swell.

\begin{cousinweb}
\textbf{\te{బొప్పి}} (\emph{boppi}) is native Dravidian: DEDR 4469 lists Telugu \emph{boppi}, \textquotedblleft{}a swelling\textquotedblright{}, in the article of \emph{poṅgu}, \textquotedblleft{}to bubble up, to swell\textquotedblright{}.
\end{cousinweb}

\begin{grammarlens}[title={What you've built}]
One more word for describing what you lived and felt.
\end{grammarlens}

\subsection*{Your turn}
\begin{itemize}
\raggedright
  \item \emph{Say it:} \emph{boppi}
  \item \emph{Say it:} \emph{boppi}, once more
  \item \emph{Say it:} \emph{talupuku tagili tammuḍi nuduṭi mīda boppi kaṭṭindi.} — My little brother hit the door and a bump came up on his forehead.
  \item \emph{Recall:} say the Telugu for joint pains, then the Telugu for to get hurt, then say \emph{boppi} again
\end{itemize}

\begin{tcolorbox}[breakable,colback=teal!4,colframe=teal!35!black,title={Before you move on}]
What does \te{బొప్పి} mean? (\textquotedblleft{}A bump; a lump on the head\textquotedblright{}.) Say it once more.
\end{tcolorbox}
\section[\texorpdfstring{\te{తూలు} (tūlu)}{తూలు (tūlu)}]{\te{తూలు} (tūlu) — to stagger; to sway}

\label{lesson:TE-C518-tulu}



\begin{quote}
Before the new one: say the Telugu for to get hurt, then the Telugu for a bump.
\end{quote}

\subsection*{You'll want to know: \te{తూలు}}
\textbf{\te{తూలు}} — \emph{tūlu} — \textquotedblleft{}to stagger; to sway\textquotedblright{}.

\begin{quote}
\te{నిద్రమత్తులో} \te{అతను} \te{తూలుతూ} \te{నడిచాడు}.
\end{quote}

\begin{quote}
\emph{nidramattulō atanu tūlutū naḍicāḍu.}
\end{quote}

\begin{quote}
Groggy with sleep, he staggered along.
\end{quote}

Past (I …): \te{తూలాను}, \emph{tūlānu}.

Unsteady on your feet, from sleep, illness or weakness.

\begin{cousinweb}
\textbf{\te{తూలు}} (\emph{tūlu}) is native Dravidian: DEDR 3376 lists Telugu \emph{tūlu}, \textquotedblleft{}to reel, to stagger, to grow feeble\textquotedblright{}.
\end{cousinweb}

\begin{grammarlens}[title={What you've built}]
One more word for describing what you lived and felt.
\end{grammarlens}

\subsection*{Your turn}
\begin{itemize}
\raggedright
  \item \emph{Say it:} \emph{tūlu}
  \item \emph{Say it:} \emph{tūlu}, once more
  \item \emph{Say it:} \emph{nidramattulō atanu tūlutū naḍicāḍu.} — Groggy with sleep, he staggered along.
  \item \emph{Recall:} say the Telugu for to get hurt, then the Telugu for a bump, then say \emph{tūlu} again
\end{itemize}

\begin{tcolorbox}[breakable,colback=teal!4,colframe=teal!35!black,title={Before you move on}]
What does \te{తూలు} mean? (\textquotedblleft{}To stagger; to sway\textquotedblright{}.) Say it once more.
\end{tcolorbox}
\section[\texorpdfstring{\te{కుంటు} (kuṇṭu)}{కుంటు (kuṇṭu)}]{\te{కుంటు} (kuṇṭu) — to limp}

\label{lesson:TE-C518-kuntu}



\begin{quote}
Before the new one: say the Telugu for a bump, then the Telugu for to stagger.
\end{quote}

\subsection*{You'll want to know: \te{కుంటు}}
\textbf{\te{కుంటు}} — \emph{kuṇṭu} — \textquotedblleft{}to limp\textquotedblright{}.

\begin{quote}
\te{కాలికి} \te{దెబ్బ} \te{తగిలి} \te{కుక్క} \te{కుంటుతోంది}.
\end{quote}

\begin{quote}
\emph{kāliki debba tagili kukka kuṇṭutōndi.}
\end{quote}

\begin{quote}
The dog hurt its leg and is limping.
\end{quote}

Past (I …): \te{కుంటాను}, \emph{kuṇṭānu}.

\te{కుంటి}, lame; \te{కుంటుతూ} \te{నడుచు}, to walk with a limp.

\begin{cousinweb}
\textbf{\te{కుంటు}} (\emph{kuṇṭu}) is native Dravidian: DEDR 1688 lists Telugu \emph{kuṇṭu}, \textquotedblleft{}to limp, to walk lamely\textquotedblright{}, and \emph{kuṇṭi}, \textquotedblleft{}lame\textquotedblright{}, with Kota \emph{kuṇṭ-}.
\end{cousinweb}

\begin{grammarlens}[title={What you've built}]
One more word for describing what you lived and felt.
\end{grammarlens}

\subsection*{Your turn}
\begin{itemize}
\raggedright
  \item \emph{Say it:} \emph{kuṇṭu}
  \item \emph{Say it:} \emph{kuṇṭu}, once more
  \item \emph{Say it:} \emph{kāliki debba tagili kukka kuṇṭutōndi.} — The dog hurt its leg and is limping.
  \item \emph{Recall:} say the Telugu for a bump, then the Telugu for to stagger, then say \emph{kuṇṭu} again
\end{itemize}

\begin{tcolorbox}[breakable,colback=teal!4,colframe=teal!35!black,title={Before you move on}]
What does \te{కుంటు} mean? (\textquotedblleft{}To limp\textquotedblright{}.) Say it once more.
\end{tcolorbox}
\section[\texorpdfstring{\te{రక్తస్రావం} (raktasrāvaṁ)}{రక్తస్రావం (raktasrāvaṁ)}]{\te{రక్తస్రావం} (raktasrāvaṁ) — bleeding; loss of blood}

\label{lesson:TE-C518-raktasravam}



\begin{quote}
Before the new one: say the Telugu for to stagger, then the Telugu for to limp.
\end{quote}

\subsection*{You'll want to know: \te{రక్తస్రావం}}
\textbf{\te{రక్తస్రావం}} — \emph{raktasrāvaṁ} — \textquotedblleft{}bleeding; loss of blood\textquotedblright{}.

\begin{quote}
\te{గాయం} \te{నుంచి} \te{రక్తస్రావం} \te{ఆగలేదు}.
\end{quote}

\begin{quote}
\emph{gāyaṁ nuñci raktasrāvaṁ āgalēdu.}
\end{quote}

\begin{quote}
The bleeding from the wound did not stop.
\end{quote}

\te{రక్తం}, blood, + \te{స్రావం}, flow.

\begin{cousinweb}
\textbf{\te{రక్తస్రావం}} (\emph{raktasrāvaṁ}) is a learned compound built from Sanskrit \emph{rakta}, \textquotedblleft{}red; blood\textquotedblright{} (CDIAL 10539), and \emph{srāva}, \textquotedblleft{}a flow\textquotedblright{} (Monier-Williams).
\end{cousinweb}

\begin{grammarlens}[title={What you've built}]
One more word for describing what you lived and felt.
\end{grammarlens}

\subsection*{Your turn}
\begin{itemize}
\raggedright
  \item \emph{Say it:} \emph{raktasrāvaṁ}
  \item \emph{Say it:} \emph{raktasrāvaṁ}, once more
  \item \emph{Say it:} \emph{gāyaṁ nuñci raktasrāvaṁ āgalēdu.} — The bleeding from the wound did not stop.
  \item \emph{Recall:} say the Telugu for to stagger, then the Telugu for to limp, then say \emph{raktasrāvaṁ} again
\end{itemize}

\begin{tcolorbox}[breakable,colback=teal!4,colframe=teal!35!black,title={Before you move on}]
What does \te{రక్తస్రావం} mean? (\textquotedblleft{}Bleeding; loss of blood\textquotedblright{}.) Say it once more.
\end{tcolorbox}
